% ---------------------------------------------------------------------
%  03 Related Work
%  Each subsection: what the line of work does, then one sentence on how
%  we differ. Do not write a literature list -- write an argument.
% ---------------------------------------------------------------------
\section{Related Work}
\label{sec:related}

\subsection{Ball detection and tracking in sports video}
\label{subsec:related_ball}
\Todo{TrackNet-style heatmap detectors, small-object detection, motion blur.
Sub-pixel localisation. How we differ: we do not stop at 2D --- 2D detections
are an intermediate representation consumed by \BLCS{}.}

\subsection{Court registration and sports field homography}
\label{subsec:related_court}
\Todo{Keypoint-based and segmentation-based field registration. Our court model
uses $\numcourtkp = 14$ ground-plane keypoints; state why this set.}

\subsection{2D and 3D human pose estimation}
\label{subsec:related_pose}
\Todo{Top-down 2D pose; 2D-to-3D lifting; world-grounded recovery
(\GVHMR{} and its predecessors). We adopt \GVHMR{} for articulated body pose
and solve the complementary problem it leaves open: placing the body in a
known metric court frame.}

\subsection{Self-supervised visual representations}
\label{subsec:related_ssl}
\Todo{\DINO{} and self-supervised pretraining for dense prediction; whether
domain-specific SSL pretraining helps small-object and thin-structure tasks.}

\subsection{3D reconstruction of sports scenes}
\label{subsec:related_3d_sports}
\Todo{Multi-view and monocular sports reconstruction; ball trajectory fitting
with physics priors. The gap this paper closes: a single reconstruction
containing both the ball and articulated players in one metric frame.}
